\subsection{Cumulative Distribution Function (CDF)}
\hrulefill

\paragraph{Definition:} The cumulative distribution function (CDF) of a continuous random variable $X$ is a function 
$F(x)$ defined as:
\[
    F(x) = P(X \leq x) = \int_{-\infty}^{x} f(t) dt
\]
where $f(t)$ is the PDF of $X$.

\paragraph{Relationship between PDF and CDF:}
\begin{itemize}
    \item The CDF $F(x)$ is the integral of the PDF $f(x)$.
    \item The PDF $f(x)$ is the derivative of the CDF $F(x)$, i.e., $f(x) = \frac{d}{dx} F(x)$.
\end{itemize}

\paragraph{Properties of CDF:}
\begin{itemize}
    \item $F(x)$ is non-decreasing.
    \item $\lim_{x \to -\infty} F(x) = 0$ and $\lim_{x \to \infty} F(x) = 1$.
    \item $P(X \leq x) = F(x)$
    \item $P(X > x) = 1 - F(x)$
    \item $P(a < X \leq b) = F(b) - F(a)$
\end{itemize}

\paragraph*{Example:}
Let $X$ be a continuous random variable with PDF:
\[
    f(x) = \begin{cases}
        \frac{2}{9} x & 0 \leq x \leq 3 \\
        0 & \text{otherwise}
    \end{cases}
\]
a. Find the CDF $F(x)$.\\
b. Use the CDF to compute $P(X > 2)$.\\

\paragraph*{a.}
\begin{align*}
    F(x) &= \int_0^x f(t) dt = \int_0^x \frac{2}{9} t dt = \begin{cases}
        0 & x < 0 \\
        \frac{2}{9} \cdot \frac{x^2}{2} = \frac{x^2}{9} & 0 \leq x \leq 3 \\
        1 & x > 3
    \end{cases}
\end{align*}

\paragraph*{b.}
\begin{align*}
    P(X > 2) &= 1 - P(X \leq 2) = 1 - F(2) \\
    &= 1 - \frac{2^2}{9} = 1 - \frac{4}{9} = \frac{5}{9}
\end{align*}

\subsubsection*{Percentiles}
\paragraph*{Definition:} The score $s$ is the $p$-th percentile $\eta (p)$ of a continuous random variable $X$ if:
\begin{itemize}
    \item $P(X \leq s) = F(s) = p$
    \item $P(X > s) = 1 - F(s) = 1 - p$
    \item $\int_{-\infty}^{\eta (p)} f(x) dx = p$
    \item $F(\eta (p)) = p$
\end{itemize}

\paragraph*{Percentile Vocabulary:}
\begin{itemize}
    \item $\eta (0.5)$ is the median.
    \item $\eta (0.25)$ is the first quartile $Q_1$.
    \item $\eta (0.75)$ is the third quartile $Q_3$.
    \item The interquartile range (IQR) is $Q_3 - Q_1$.
\end{itemize}

\paragraph*{Example:}
Find the 75th percentile of the continuous random variable $X$ from the previous example.
\begin{align*}
    \eta (0.75) &= s \text{ such that } F(s) = 0.75 \\
    F(s) &= \frac{s^2}{9} = 0.75 \\
    s^2 &= 6.75 \\
    s &= \sqrt{6.75} \approx 2.598
\end{align*}

\hrulefill

\subsubsection{Mean and Variance of a Continuous Random Variable}
\paragraph{Mean (Expected Value):} The mean or expected value of a continuous random variable $X$ with PDF $f(x)$ is given by:
\begin{align*}
    E(X) &= \int_{-\infty}^{\infty} x f(x) dx\\
    E(h(X)) &= \int_{-\infty}^{\infty} h(x) f(x) dx\\
    E(aX+b) &= aE(X) + b\\
    E(h(X)) &\neq h(E(X)) \text{ unless $h(x)$ is linear}
\end{align*}


\paragraph{Variance:} The variance of a continuous random variable $X$ with PDF $f(x)$ is given by:
\begin{align*}
    \sigma^2 &= Var(X) = \int_{-\infty}^{\infty} (x - \mu)^2 f(x) dx \\
    Var(X) &= E(X^2) - (E(X))^2\\
    Var(aX+b) &= a^2 Var(X)
\end{align*}

\pagebreak

\paragraph*{Example:}
Let $X$ be a continuous random variable with PDF:
\[
    f(x) = \begin{cases}
        \frac{3}{8} x^2 & 0 \leq x \leq 2 \\
        0 & \text{otherwise}
    \end{cases}
\]
a. Find $E(X)$\\
b. Find $E(X^2)$\\
c. Find $Var(X)$\\

\paragraph*{a.}
\begin{align*}
    E(X) &= \int_0^2 x \cdot \frac{3}{8} x^2 dx = \int_0^2 \frac{3}{8} x^3 dx = \left[ \frac{3}{8} \cdot \frac{x^4}{4} \right]_0^2 = \frac{3}{8} \cdot 4 = \frac{3}{2}
\end{align*}

\paragraph*{b.}
\begin{align*}
    E(X^2) &= \int_0^2 x^2 \cdot \frac{3}{8} x^2 dx = \int_0^2 \frac{3}{8} x^4 dx = \left[ \frac{3}{8} \cdot \frac{x^5}{5} \right]_0^2 = \frac{3}{8} \cdot \frac{32}{5} = \frac{12}{5}
\end{align*}

\paragraph*{c.}
\begin{align*}
    Var(X) &= E(X^2) - (E(X))^2 = \frac{12}{5} - \left( \frac{3}{2} \right)^2 = \frac{12}{5} - \frac{9}{4} = \frac{48}{20} - \frac{45}{20} = \frac{3}{20}
\end{align*}